% page 626 of the facsimile (image p-625 of the scan)
% block 154.9 x 246.3 mm ; left margin 8.0 mm
\pgc{-12.700mm}{0.000mm}{
\l{\cel{7}618}
\l{}
\l{\sou{violentar}. (\sou{trans.})\cel{2}Koaktar sovaje, per la uzo di vigoro}
\l{\cel{3}sen-represa. - EF.}
\l{}
\l{violino. (\sou{muziko)}. Instrumento kun kordi ed arketo, derivita}
\l{\cel{3}de la viulo olima, igita plu mikra-dimensiona, e reduk-}
\l{\cel{3}tita a quar kordi inter-akordita quintope (g, d, la, e).}
\l{\cel{3}- DEFIS.}
\l{}
\l{violoncelo.\cel{2}(\sou{muziko}).Instrumento kun quar kordi ed arketo,}
\l{\cel{3}ye un quinto sub alto, qua preske korespondas a la viulo-}
\l{\cel{3}baso olima, e quan onu pleas tenante lu inter sua genui.}
\l{\cel{3}- DEFIS.}
\l{}
\l{\sou{vipero}. (\sou{zool.}) Reptero serpenta di qua la morduro esas ve-}
\l{\cel{3}nenoza. L.\cel{1}\sou{vipera}. - DEFIS.}
\l{}
\l{\sou{viro}. Adultulo homa, qua havas maxim-alta-grade la karakter-}
\l{\cel{3}izivi di la maskuli. - EFIS.}
\l{}
\l{\sou{virga}. I. (\sou{pri adolecantino homa}). Quan nula viro koitis. -}
\l{\cel{3}II. (metaf.) (a) Regiono virga : ne ja explorita ; (b)}
\l{\cel{3}Tero virga : ne ja kultivita ; (c) Foresto virga : neja}
\l{\cel{3}explotita ; (d)\cel{1}\sou{Metalo virga}\cel{1}: renkontrita pura en sulo.}
\l{\cel{3}- EFIS.}
\l{}
\l{\sou{virtuala}. I. (\sou{filoz.}) Qua povas efikar, ma ne ja agis lo.}
\l{\cel{3}II. (optiko) : La foko virtuala di spegulo : loko ube la}
\l{\cel{3}lumo-radii diverganta quin reflektas la spegulo, se}
\l{\cel{3}prolongite ideale, konvergus dop la spegulo. - III. (me-}
\l{\cel{3}\sou{kaniko})\cel{3}\sou{Efiko virtuala}\cel{1}: Efiko quan agus moveblo se lu}
\l{\cel{3}movesus ye quanto extreme mikra ye ta o ca instanto. -}
\l{\cel{3}DEFIRS.}
\l{}
\l{\sou{virtuoza.}\cel{1}Di qua la talento igas (ta persono) ecepto-kazo.}
\l{\cel{3}DEFIS.}
\l{}
\l{\sou{virulenta}. (\sou{patol}.) Qua kontenas la fonto di infekto e di}
\l{\cel{3}transmiso morbala. - DEFIS.}
\l{}
\l{\sou{viruso}. (\sou{patol.}) Substanco qua esas la kauzo di infekto e}
\l{\cel{3}di transmiso morbala. - DEFIS.}
\l{}
\l{\sou{visar}. (\sou{trans.})\cel{2}(\sou{tekn.}) Rotacigar skrubo por sinkar lu (ad-}
\l{\cel{3}en ulo). - F.}
\l{}
\l{\sou{vishar}. (\sou{trans.}) Netigar (objekto) per pasigar sur olu ulo}
\l{\cel{3}qua deprenas la strateto de polvo qua kovras lu; sekigar}
\l{\cel{3}per froto. - D.}
\l{}
\l{\sou{viskomto}. I. Titulo siniorala, aplikata,en la komenco, a}
\l{\cel{3}la lietnanto di la komto. - II. Nobeleso-titulo, sub komto}
\l{\cel{3}e super barono. - DEFIRS.}
\l{}
\l{\sou{viskoza}. Di qua la molekuli esas tenaca, inter-adheras ed}
\l{\cel{3}adheras a la korpo quan li tushas. - DEFIS.}
}
